% !TeX root = closed-systems-from-comparison-completeness.tex
% ============================================================
\chapter{Observer-Level Irreversible Descent and the Two-Locus
Residue}
\label{ch:flow}
% ============================================================
\section{Introduction}
Under the standing principle of closed-world admissibility
(\cref{sp:closed-world}), this chapter sets up the internal
observer-reduction framework at fixed quotient semantics, i.e.
relative to \cref{sp:quotient,sp:transport}, and isolates the
conditional regime in which observer-level irreversibility can
arise.
Using the static transport obstruction of \cref{ch:bridge}, it sets
up the internal observer-reduction framework and the conditional
irreversibility criterion carried next into \cref{ch:stitching};
only later, together with the stabilized interface-side package, do
those observer-level structures feed \cref{ch:causality}.
\Cref{ch:bottom,ch:rotation,ch:lift,ch:descent_obstruction,ch:bridge}
established two structural facts.
First, closed-system comparison semantics determines a canonical diagonal
quotient
\[
	\pi\colon X\to\Phys:=X/G,
\]
which removes precisely the redundant degrees of freedom coming from the
intrinsic diagonal symmetry. This quotient is internal to the closed
system. It is not an observer coarse-graining, and by itself it creates
no arrow of time.
Second, any enrichment of quotient-level protocol data back to
representatives in \(X\) remains exhausted by the same two previously
classified loci, possibly both and with no third enrichment locus:
representative choice and morphism-level transport or holonomy data.
The purpose of the present chapter is to determine how this quotient
architecture separates intrinsic reversible descent from the later
conditional observer-induced descent regime.
There are two logically distinct descent stages:
\[
	X\xrightarrow{\ \pi\ }\Phys\xrightarrow{\ \theta\ }Y.
\]
\begin{enumerate}[label=\textup{(\roman*)}]
	\item
	      The intrinsic quotient
	      \[
		      \pi\colon X\to\Phys
	      \]
	      preserves reversibility: a reversible \(G\)-equivariant flow on \(X\)
	      descends to a reversible flow on \(\Phys\).
	\item
	      Irreversibility may arise only after a further internal observer
	      reduction
	      \[
		      \theta\colon \Phys\twoheadrightarrow Y,
	      \]
	      built from restricted admissible quotient-level observations;
	      the chapter's irreversibility criterion then applies only once a
	      proper such observer family is already forward-compatible with
	      the descended reversible flow on \(\Phys\).
\end{enumerate}
The first task of the chapter is to isolate this two-stage architecture
and, once a compatible irreversible descent is fixed, derive the
resulting filtration of indistinguishability relations on \(\Phys\).
From that filtration follow two canonical consequences:
a monotone entropy functional and an extended pseudo-ultrametric
geometry.
The second task is structural. We prove that irreversibility at the
observer level introduces no new enrichment primitive. Even after the
reduction \(\theta\), every compositional lift of quotient-level data
back to representatives in \(X\) remains exhausted by the same two loci
identified in
\cref{thm:lift-classification-main,cor:two-loci-lift}: representative
choice and morphism-level transport or holonomy data, possibly both,
with no third enrichment locus.
The central structural claim is therefore the following: once a
\(G\)-equivariant reversible flow on \(X\) and its descended reversible
flow on \(\Phys\) are fixed, nontrivial transport obstruction forces
irreversible observer descent whenever a proper internal observer
family built from restricted quotient-level observations is already
forward-compatible with that descended flow.
\medskip
\noindent
\textbf{Dependence on labeled results.}
This chapter uses:
\begin{itemize}
	\item the intrinsic quotient semantics \(X\to\Phys\) established in the
	      quotient chapters;
	\item representative lift classification and the two-locus analysis from
	      Chapters~\ref{ch:lift} and \ref{ch:descent_obstruction};
	\item the triangle--loop bridge and nontrivial transport obstruction
	      from Chapter~\ref{ch:bridge}.
\end{itemize}
\medskip
\noindent
\textbf{Use in subsequent chapters.}
This chapter provides:
\begin{itemize}
	\item the rigid separation between intrinsic reversible quotient
	      semantics and the later observer-level descent on \(\Phys\) where
	      irreversible behavior may arise;
	\item the internal observer-reduction framework on \(\Phys\),
	      together with the conditional irreversibility criterion for
	      proper observer families built from restricted quotient-level
	      protocols once forward compatibility is fixed;
	\item the canonical filtration of observer indistinguishability on
	      \(\Phys\) attached to a compatible irreversible descent;
	\item the entropy and extended pseudo-ultrametric structures induced
	      by that filtration;
	\item the theorem that even in the irreversible regime the
	      enrichment residue above quotient semantics remains exhausted by
	      the same two loci: representative choice and morphism-level
	      transport or holonomy data, possibly both, with no third locus.
\end{itemize}
% ------------------------------------------------------------
\section{Reversible flows and forward factors}
% ------------------------------------------------------------
We begin with the abstract distinction between reversible dynamics and
forward-only reduced dynamics.
\begin{definition}[Reversible flow]
	A \emph{reversible flow} on a set \(X\) is a family
	\[
		(\Phi_t)_{t\in\mathbb R}
	\]
	of bijections \(\Phi_t\colon X\to X\) satisfying
	\[
		\Phi_0=\id_X,
		\qquad
		\Phi_{t+s}=\Phi_t\circ\Phi_s
		\qquad
		\text{for all }s,t\in\mathbb R.
	\]
\end{definition}
\begin{definition}[Forward factor map]
	Let \((\Phi_t)_{t\in\mathbb R}\) be a reversible flow on \(X\).
	A surjection
	\[
		\rho\colon X\twoheadrightarrow Y
	\]
	is a \emph{forward factor map} if for every \(t\ge0\) there exists a map
	\[
		\Psi_t\colon Y\to Y
	\]
	such that
	\[
		\rho\circ\Phi_t=\Psi_t\circ\rho.
	\]
\end{definition}
\begin{lemma}[Forward factor semigroup]
	If \(\rho\colon X\twoheadrightarrow Y\) is a forward factor map, then
	the family \((\Psi_t)_{t\ge0}\) satisfies
	\[
		\Psi_0=\id_Y,
		\qquad
		\Psi_{t+s}=\Psi_t\circ\Psi_s
		\qquad
		\text{for all }s,t\ge0.
	\]
\end{lemma}
\begin{proof}
	Since
	\[
		\rho\circ\Phi_0=\rho,
	\]
	surjectivity of \(\rho\) implies
	\[
		\Psi_0=\id_Y.
	\]
	Let \(s,t\ge0\). Then
	\[
		\Psi_t\circ\Psi_s\circ\rho
		=
		\Psi_t\circ\rho\circ\Phi_s
		=
		\rho\circ\Phi_t\circ\Phi_s
		=
		\rho\circ\Phi_{t+s}
		=
		\Psi_{t+s}\circ\rho.
	\]
	Again by surjectivity of \(\rho\),
	\[
		\Psi_{t+s}=\Psi_t\circ\Psi_s.
	\]
\end{proof}
\begin{definition}[Information loss]
	\label{def:flow-information-loss}
	A forward factor map \(\rho\) \emph{loses information} if there exist
	distinct points \(x_1,x_2\in X\) and a time \(t_0>0\) such that
	\[
		\rho(x_1)=\rho(x_2)
	\]
	but
	\[
		\rho(\Phi_{-t_0}(x_1))\neq \rho(\Phi_{-t_0}(x_2)).
	\]
\end{definition}
Thus \(\rho\) identifies states which remain distinguishable when the
underlying reversible dynamics is run backward.
\begin{lemma}[Information loss obstructs two-sided descent]
	\label{lem:flow-information-loss}
	If \(\rho\) loses information, then there does not exist a family
	\[
		(\widetilde\Psi_t)_{t\in\mathbb R}
	\]
	extending \((\Psi_t)_{t\ge0}\) and satisfying
	\[
		\rho\circ\Phi_t=\widetilde\Psi_t\circ\rho
		\qquad
		\text{for all }t\in\mathbb R.
	\]
\end{lemma}
\begin{proof}
	Assume for contradiction that such an extension exists.
	Choose \(x_1,x_2,t_0\) witnessing information loss. Evaluating the
	intertwining identity at time \(-t_0\) gives
	\[
		\widetilde\Psi_{-t_0}(\rho(x_1))
		=
		\rho(\Phi_{-t_0}(x_1)),
		\qquad
		\widetilde\Psi_{-t_0}(\rho(x_2))
		=
		\rho(\Phi_{-t_0}(x_2)).
	\]
	Since \(\rho(x_1)=\rho(x_2)\), the two left-hand sides are equal,
	whereas the two right-hand sides are distinct by hypothesis. This is
	impossible.
\end{proof}
\begin{remark}[Irreversibility mechanism]
	Noninjectivity of \(\rho\) alone does not force irreversibility.
	The obstruction is the backward fiber separation of
	\cref{def:flow-information-loss}. Irreversibility appears exactly when
	states that are identified at the reduced level fail to remain
	identified under backward evolution in the unreduced reversible system.
\end{remark}
% ------------------------------------------------------------
\section{The intrinsic quotient preserves reversibility}
% ------------------------------------------------------------
We now apply the abstract discussion to the closed-system quotient
architecture.
\begin{definition}[\(G\)-equivariant flow]
	A reversible flow \((\Phi_t)_{t\in\mathbb R}\) on \(X\) is
	\emph{\(G\)-equivariant} if
	\[
		\Phi_t(g\cdot x)=g\cdot\Phi_t(x)
		\qquad
		\text{for all }g\in G,\ x\in X,\ t\in\mathbb R.
	\]
\end{definition}
\begin{proposition}[Time descends through the intrinsic quotient]
	\label{prop:flow-time-descends}
	Suppose \((\Phi_t)_{t\in\mathbb R}\) is a \(G\)-equivariant reversible
	flow on \(X\). Then there exists a unique reversible flow
	\[
		(\widehat\Psi_t)_{t\in\mathbb R}
	\]
	on \(\Phys\) such that
	\[
		\pi\circ\Phi_t=\widehat\Psi_t\circ\pi
		\qquad
		\text{for all }t\in\mathbb R.
	\]
\end{proposition}
\begin{proof}
	Define
	\[
		\widehat\Psi_t([x]):=[\Phi_t(x)].
	\]
	To prove well-definedness, suppose \([x]=[y]\). Then \(y=g\cdot x\) for
	some \(g\in G\). By equivariance,
	\[
		\Phi_t(y)=\Phi_t(g\cdot x)=g\cdot\Phi_t(x),
	\]
	so \([\Phi_t(y)]=[\Phi_t(x)]\). Thus \(\widehat\Psi_t\) is well defined.
	For every \(x\in X\),
	\[
		(\widehat\Psi_t\circ\pi)(x)=\widehat\Psi_t([x])=[\Phi_t(x)]
		=
		(\pi\circ\Phi_t)(x),
	\]
	so the intertwining identity holds.
	The flow laws follow directly from those of \((\Phi_t)\):
	\[
		\widehat\Psi_0([x])=[x],
	\]
	and
	\[
		\widehat\Psi_{t+s}([x])
		=
		[\Phi_{t+s}(x)]
		=
		[\Phi_t(\Phi_s(x))]
		=
		\widehat\Psi_t(\widehat\Psi_s([x])).
	\]
	Since each \(\Phi_t\) is bijective, each \(\widehat\Psi_t\) is bijective
	with inverse \(\widehat\Psi_{-t}\). Hence
	\((\widehat\Psi_t)_{t\in\mathbb R}\) is a reversible flow on \(\Phys\).
	Uniqueness follows from surjectivity of \(\pi\).
\end{proof}
\begin{corollary}[No arrow from the intrinsic quotient]
	If
	\[
		\pi(x_1)=\pi(x_2),
	\]
	then for every \(t\ge0\) one has
	\[
		\pi(\Phi_{-t}(x_1))=\pi(\Phi_{-t}(x_2)).
	\]
\end{corollary}
\begin{proof}
	By \cref{prop:flow-time-descends}, the descended reversible flow
	\((\widehat\Psi_t)_{t\in\mathbb R}\) exists on \(\Phys\). Therefore
	\[
		\pi(\Phi_{-t}(x_1))
		=
		\widehat\Psi_{-t}(\pi(x_1))
		=
		\widehat\Psi_{-t}(\pi(x_2))
		=
		\pi(\Phi_{-t}(x_2)).
	\]
\end{proof}
Thus the intrinsic diagonal quotient preserves reversibility.
Any arrow of time must therefore arise only after a later internal
observer reduction
\[
	\theta\colon \Phys\twoheadrightarrow Y.
\]
% ------------------------------------------------------------
\section{Internal observer reductions}
% ------------------------------------------------------------
The second descent stage must be internal to the closed system.
Accordingly, we now set up the internal observer-reduction framework
using restricted quotient-level observation protocols.
\begin{definition}[Admissible quotient-level observable]
	An \emph{admissible quotient-level observable} is a map
	\[
		\omega\colon \Phys\to V
	\]
	which arises from an admissible quotient-level protocol functor and is
	invariant under the transport equivalences already determined at the
	quotient level.
\end{definition}
\begin{definition}[Observable family]
	A \emph{quotient-level observable family} is a collection
	\[
		\mathfrak O=\{\omega_a\colon \Phys\to V_a\}_{a\in A}
	\]
	of admissible quotient-level observables.
\end{definition}
\begin{definition}[Observer indistinguishability]
	Given a quotient-level observable family \(\mathfrak O\), define a
	relation \(\sim_{\mathfrak O}\) on \(\Phys\) by
	\[
		p\sim_{\mathfrak O}q
		\iff
		\omega_a(p)=\omega_a(q)
		\qquad
		\text{for all }a\in A.
	\]
\end{definition}
\begin{lemma}[Observer indistinguishability is an equivalence relation]
	For every quotient-level observable family \(\mathfrak O\), the relation
	\(\sim_{\mathfrak O}\) is an equivalence relation on \(\Phys\).
\end{lemma}
\begin{proof}
	Reflexivity and symmetry are immediate. For transitivity, suppose
	\[
		p\sim_{\mathfrak O}q
		\qquad
		\text{and}
		\qquad
		q\sim_{\mathfrak O}r.
	\]
	Then for every \(a\in A\),
	\[
		\omega_a(p)=\omega_a(q)=\omega_a(r).
	\]
	Hence \(p\sim_{\mathfrak O}r\).
\end{proof}
\begin{definition}[Internal observer reduction]
	Let \(\mathfrak O\) be a quotient-level observable family. Define
	\[
		Y_{\mathfrak O}:=\Phys/{\sim_{\mathfrak O}},
	\]
	and let
	\[
		\theta_{\mathfrak O}\colon \Phys\twoheadrightarrow Y_{\mathfrak O}
	\]
	be the canonical quotient map
	\[
		\theta_{\mathfrak O}(p):=[p]_{\sim_{\mathfrak O}}.
	\]
	This map is called the \emph{internal observer reduction} induced by
	\(\mathfrak O\).
\end{definition}
\begin{remark}[Internality]
	The map \(\theta_{\mathfrak O}\) is not an externally imposed
	coarse-graining. It is determined entirely by a restricted family of
	admissible quotient-level observations already internal to the closed
	system.
\end{remark}
\begin{proposition}[Universal property of observer reduction]
	For every \(\omega_a\in\mathfrak O\), there exists a unique map
	\[
		\overline\omega_a\colon Y_{\mathfrak O}\to V_a
	\]
	such that
	\[
		\omega_a=\overline\omega_a\circ\theta_{\mathfrak O}.
	\]
	Moreover, \(\theta_{\mathfrak O}\) is initial among surjections with
	this property.
\end{proposition}
\begin{proof}
	If \(p\sim_{\mathfrak O}q\), then \(\omega_a(p)=\omega_a(q)\), so
	\(\omega_a\) is constant on \(\sim_{\mathfrak O}\)-classes. Therefore
	\[
		\overline\omega_a([p]_{\sim_{\mathfrak O}}):=\omega_a(p)
	\]
	is well defined. Uniqueness follows from surjectivity of
	\(\theta_{\mathfrak O}\).
	Now suppose
	\[
		\eta\colon \Phys\twoheadrightarrow Z
	\]
	is another surjection through which every \(\omega_a\) factors. If
	\(\eta(p)=\eta(q)\), then \(\omega_a(p)=\omega_a(q)\) for every \(a\),
	hence \(p\sim_{\mathfrak O}q\). Therefore \(\eta\) refines
	\(\sim_{\mathfrak O}\), and there exists a unique map
	\[
		u\colon Z\to Y_{\mathfrak O}
	\]
	with
	\[
		\theta_{\mathfrak O}=u\circ\eta.
	\]
	Thus \(\theta_{\mathfrak O}\) is initial among such surjections.
\end{proof}
\begin{definition}[Proper observer family]
	A quotient-level observable family \(\mathfrak O\) is \emph{proper} if
	there exist distinct points \(p_1,p_2\in \Phys\) such that
	\[
		p_1\sim_{\mathfrak O}p_2.
	\]
	Equivalently, the induced observer reduction
	\(\theta_{\mathfrak O}\) is noninjective.
\end{definition}
\begin{definition}[Forward-compatible observer family]
	Let \((\widehat\Psi_t)_{t\in\mathbb R}\) be a reversible flow on
	\(\Phys\). A quotient-level observable family \(\mathfrak O\) is
	\emph{forward-compatible} with \((\widehat\Psi_t)\) if for every
	\(t\ge0\) there exists a map
	\[
		\Psi_t^{\mathfrak O}\colon Y_{\mathfrak O}\to Y_{\mathfrak O}
	\]
	such that
	\[
		\theta_{\mathfrak O}\circ\widehat\Psi_t
		=
		\Psi_t^{\mathfrak O}\circ\theta_{\mathfrak O}.
	\]
\end{definition}
\begin{definition}[Backward-stable observer family]
	Let \((\widehat\Psi_t)_{t\in\mathbb R}\) be a reversible flow on
	\(\Phys\), and let \(\mathfrak O\) be a quotient-level observable
	family that is forward-compatible with \((\widehat\Psi_t)\).
	We say that \(\mathfrak O\) is \emph{backward-stable} if the
	intertwining identity extends to all \(t\in\mathbb R\), i.e. if there
	exists a family
	\[
		(\widetilde\Psi_t^{\mathfrak O})_{t\in\mathbb R}
	\]
	on \(Y_{\mathfrak O}\) such that
	\[
		\theta_{\mathfrak O}\circ\widehat\Psi_t
		=
		\widetilde\Psi_t^{\mathfrak O}\circ\theta_{\mathfrak O}
		\qquad
		\text{for all }t\in\mathbb R.
	\]
\end{definition}
% ------------------------------------------------------------
\section{Transport obstruction and fixed-flow observer irreversibility}
% ------------------------------------------------------------
We now identify the exact point at which, relative to fixed reversible
flow data, nontrivial transport obstruction forces observer-level
irreversibility once the associated restricted observer family is
forward-compatible.
\begin{definition}[Observer-level information loss on \(\Phys\)]
	\label{def:flow-observer-information-loss}
	Let \((\widehat\Psi_t)_{t\in\mathbb R}\) be a reversible flow on
	\(\Phys\), and let \(\mathfrak O\) be a quotient-level observable
	family that is forward-compatible with \((\widehat\Psi_t)\).
	We say that \(\mathfrak O\) \emph{loses information} if there exist
	distinct points \(p_1,p_2\in \Phys\) and a time \(t_0>0\) such that
	\[
		\theta_{\mathfrak O}(p_1)=\theta_{\mathfrak O}(p_2)
	\]
	but
	\[
		\theta_{\mathfrak O}(\widehat\Psi_{-t_0}(p_1))
		\neq
		\theta_{\mathfrak O}(\widehat\Psi_{-t_0}(p_2)).
	\]
\end{definition}
\begin{lemma}[Lifted and quotient-level information loss agree]
	\label{lem:flow-information-loss-compare}
	Let \((\Phi_t)\) be a \(G\)-equivariant reversible flow on \(X\), let
	\((\widehat\Psi_t)\) be the descended reversible flow on \(\Phys\), and
	let \(\mathfrak O\) be a quotient-level observable family on
	\(\Phys\) that is forward-compatible with \((\widehat\Psi_t)\). Set
	\[
		\theta:=\theta_{\mathfrak O},
		\qquad
		\rho:=\theta\circ\pi.
	\]
	Then \(\rho\) loses information in the sense of
	\cref{def:flow-information-loss} if and only if \(\mathfrak O\) loses
	information in the sense of \cref{def:flow-observer-information-loss}.
\end{lemma}
\begin{proof}
	Assume first that \(\rho\) loses information. Then there exist
	\(x_1,x_2\in X\) and \(t_0>0\) such that
	\[
		\rho(x_1)=\rho(x_2)
	\]
	but
	\[
		\rho(\Phi_{-t_0}(x_1))\neq \rho(\Phi_{-t_0}(x_2)).
	\]
	Set
	\[
		p_i:=\pi(x_i)\in\Phys
		\qquad
		(i=1,2).
	\]
	Then
	\[
		\theta(p_1)=\rho(x_1)=\rho(x_2)=\theta(p_2).
	\]
	Moreover, by \cref{prop:flow-time-descends},
	\[
		\pi(\Phi_{-t_0}(x_i))=\widehat\Psi_{-t_0}(p_i).
	\]
	Hence
	\[
		\theta(\widehat\Psi_{-t_0}(p_1))
		=
		\rho(\Phi_{-t_0}(x_1))
		\neq
		\rho(\Phi_{-t_0}(x_2))
		=
		\theta(\widehat\Psi_{-t_0}(p_2)).
	\]
	Thus \(\mathfrak O\) loses information on \(\Phys\).
	Conversely, assume \(\mathfrak O\) loses information on \(\Phys\). Then
	there exist distinct points \(p_1,p_2\in \Phys\) and \(t_0>0\) such that
	\[
		\theta(p_1)=\theta(p_2)
	\]
	but
	\[
		\theta(\widehat\Psi_{-t_0}(p_1))
		\neq
		\theta(\widehat\Psi_{-t_0}(p_2)).
	\]
	Choose representatives \(x_i\in X\) with
	\[
		\pi(x_i)=p_i
		\qquad
		(i=1,2).
	\]
	Then
	\[
		\rho(x_1)=\theta(\pi(x_1))=\theta(p_1)=\theta(p_2)=\theta(\pi(x_2))=\rho(x_2),
	\]
	while again by \cref{prop:flow-time-descends},
	\[
		\rho(\Phi_{-t_0}(x_i))
		=
		\theta(\pi(\Phi_{-t_0}(x_i)))
		=
		\theta(\widehat\Psi_{-t_0}(p_i)).
	\]
	Hence
	\[
		\rho(\Phi_{-t_0}(x_1))
		\neq
		\rho(\Phi_{-t_0}(x_2)).
	\]
	Thus \(\rho\) loses information.
\end{proof}
\begin{definition}[Compatible irreversible descent]
	A \emph{compatible irreversible descent} consists of:
	\begin{enumerate}[label=\textup{(\arabic*)}]
		\item a \(G\)-equivariant reversible flow
		      \[
			      (\Phi_t)_{t\in\mathbb R}
		      \]
		      on \(X\);
		\item a proper quotient-level observable family \(\mathfrak O\) on
		      \(\Phys\), forward-compatible with the descended reversible flow
		      \((\widehat\Psi_t)_{t\in\mathbb R}\) on \(\Phys\) determined by
		      item \textup{(1)} via \cref{prop:flow-time-descends};
		\item the induced internal observer reduction
		      \[
			      \theta:=\theta_{\mathfrak O}\colon \Phys\twoheadrightarrow Y_{\mathfrak O};
		      \]
		\item information loss for \(\mathfrak O\), equivalently for
		      \(\rho:=\theta\circ\pi\), in the sense of
		      \cref{lem:flow-information-loss-compare}.
	\end{enumerate}
\end{definition}
\begin{lemma}[Factor identity on \(\Phys\)]
	\label{lem:flow-factor-identity}
	Let
	\[
		\rho=\theta\circ\pi\colon X\twoheadrightarrow Y
	\]
	be a compatible irreversible descent, let
	\((\widehat\Psi_t)_{t\in\mathbb R}\) be the descended reversible flow on
	\(\Phys\), and let
	\((\Psi_t)_{t\ge0}\) be the induced forward semigroup on \(Y\).
	Then for every \(t\ge0\),
	\[
		\theta\circ\widehat\Psi_t=\Psi_t\circ\theta.
	\]
\end{lemma}
\begin{proof}
	For \(x\in X\) and \(t\ge0\),
	\[
		(\theta\circ\widehat\Psi_t\circ\pi)(x)
		=
		(\theta\circ\pi\circ\Phi_t)(x)
		=
		(\rho\circ\Phi_t)(x)
		=
		(\Psi_t\circ\rho)(x)
		=
		(\Psi_t\circ\theta\circ\pi)(x).
	\]
	Since \(\pi\) is surjective,
	\[
		\theta\circ\widehat\Psi_t=\Psi_t\circ\theta.
	\]
\end{proof}
We now state the supported conditional forcing theorem.
\begin{theorem}[Conditional fixed-flow observer irreversibility]
	\label{thm:flow-inevitable-loss}
	Under the standing principle \cref{sp:closed-world}, assume the closed
	comparison system is nontrivial in the sense that the
	transport obstruction of Chapter~\ref{ch:bridge} is nonvanishing on some
	finite transport subsystem. Let \((\Phi_t)_{t\in\mathbb R}\) be a
	\(G\)-equivariant reversible flow on \(X\), and let
	\((\widehat\Psi_t)_{t\in\mathbb R}\) be the descended reversible flow on
	\(\Phys\). Fix a minimal triangle subsystem on which transport fails to
	descend through triangle coherence, and let \(\mathfrak O\) be the
	proper admissible quotient-level observable family on \(\Phys\) whose
	values depend only on forward-coherent restrictions of that subsystem.
	Assume that \(\mathfrak O\) is forward-compatible with
	\((\widehat\Psi_t)\). Then:
	\begin{enumerate}[label=\textup{(\roman*)}]
		\item \(\mathfrak O\) is not backward-stable with respect to
		      \((\widehat\Psi_t)\);
		\item the induced internal observer reduction
		      \[
			      \theta_{\mathfrak O}\colon \Phys\twoheadrightarrow Y_{\mathfrak O}
		      \]
		      loses information in the sense of
		      \cref{def:flow-observer-information-loss};
		\item the composite
		      \[
			      \rho:=\theta_{\mathfrak O}\circ\pi
		      \]
		      loses information on \(X\), hence defines a compatible
		      irreversible descent.
	\end{enumerate}
\end{theorem}
\begin{proof}
	By \cref{thm:bridge-main}, nontrivial transport obstruction means that
	there exists a finite transport subsystem carrying a minimal triangle
	regime, equivalently a directed triangle subsystem on which transport
	fails to descend through triangle coherence. Fix such a subsystem, and
	let \(\mathfrak O\) be the associated admissible quotient-level
	observable family from the statement. This family is internal, because
	it is defined entirely in terms of admissible quotient-level protocols,
	and it is proper because the full triangle obstruction is invisible on
	every two-edge restriction while distinct global states remain.
	Let \(\theta_{\mathfrak O}\colon\Phys\twoheadrightarrow Y_{\mathfrak O}\)
	be the induced observer reduction. By assumption, \(\mathfrak O\) is
	forward-compatible with the fixed descended reversible flow
	\((\widehat\Psi_t)_{t\in\mathbb R}\) on \(\Phys\).
	We claim that \(\mathfrak O\) is not backward-stable. Suppose instead
	that it were backward-stable with respect to \((\widehat\Psi_t)\). Then
	there would exist a family
	\[
		(\widetilde\Psi_t^{\mathfrak O})_{t\in\mathbb R}
	\]
	on \(Y_{\mathfrak O}\) such that
	\[
		\theta_{\mathfrak O}\circ\widehat\Psi_t
		=
		\widetilde\Psi_t^{\mathfrak O}\circ\theta_{\mathfrak O}
		\qquad
		\text{for all }t\in\mathbb R.
	\]
	On the chosen minimal obstructed subsystem, the family \(\mathfrak O\)
	records only the forward-coherent restricted reports. A two-sided
	intertwining through \(\theta_{\mathfrak O}\) would therefore make those
	restricted reports stable under full transport and force the missing
	global coherence to descend through the observer quotient. By the
	triangle--loop identification of \cref{thm:bridge-main}, that would give
	a two-sided transport descent across the transport-obstructed subsystem,
	contradicting the assumed obstruction. Hence \(\mathfrak O\) is not
	backward-stable.
	If, for every \(t>0\) and every \(p_1,p_2\in\Phys\) with
	\(\theta_{\mathfrak O}(p_1)=\theta_{\mathfrak O}(p_2)\), one also had
	\[
		\theta_{\mathfrak O}(\widehat\Psi_{-t}(p_1))
		=
		\theta_{\mathfrak O}(\widehat\Psi_{-t}(p_2)),
	\]
	then each backward iterate \(\widehat\Psi_{-t}\) would descend through
	the surjection \(\theta_{\mathfrak O}\) to a map on \(Y_{\mathfrak O}\).
	Together with forward compatibility for \(t\ge0\), this would make
	\(\mathfrak O\) backward-stable, a contradiction. Therefore there exist
	distinct points \(p_1,p_2\in \Phys\) and a time \(t_0>0\) such that
	\[
		\theta_{\mathfrak O}(p_1)=\theta_{\mathfrak O}(p_2)
	\]
	but
	\[
		\theta_{\mathfrak O}(\widehat\Psi_{-t_0}(p_1))
		\neq
		\theta_{\mathfrak O}(\widehat\Psi_{-t_0}(p_2)).
	\]
	That is exactly observer-level information loss in the sense of
	\cref{def:flow-observer-information-loss}. By
	\cref{lem:flow-information-loss-compare}, the composite
	\[
		\rho:=\theta_{\mathfrak O}\circ\pi
	\]
	loses information on \(X\) as well. Since \(\mathfrak O\) is proper and
	forward-compatible with the fixed descended flow, \(\rho\) is therefore
	a compatible irreversible descent.
\end{proof}
\begin{corollary}[Conditional irreversibility for a fixed flow]
	\label{cor:flow-irreversibility-forced}
	Assume the hypotheses of \cref{thm:flow-inevitable-loss}. Then, for the
	fixed reversible flow data, the associated internal observer reduction
	\(\theta_{\mathfrak O}\) is forward-compatible but not backward-stable
	and therefore exhibits observer-level irreversibility.
\end{corollary}
\begin{proof}
	Fix the \(G\)-equivariant reversible flow \((\Phi_t)\) on \(X\), the
	descended reversible flow \((\widehat\Psi_t)\) on \(\Phys\), and the
	associated observer family \(\mathfrak O\) as in
	\cref{thm:flow-inevitable-loss}. The theorem gives non-backward-stability
	and observer-level information loss for \(\theta_{\mathfrak O}\), which
	is exactly the claimed irreversibility.
\end{proof}
\begin{remark}[Internal source of the arrow]
	The observer reduction \(\theta\) is now fully internalized: it is
	induced by a proper restricted family of admissible quotient-level
	observations. Irreversibility is therefore not postulated by an
	external coarse-graining. For the fixed reversible flow data of
	\cref{thm:flow-inevitable-loss}, it appears precisely when internal
	observer reduction is forward-compatible with the descended flow but not
	backward-stable.
\end{remark}
% ------------------------------------------------------------
\section{\texorpdfstring{Time-indexed indistinguishability on \(\Phys\)}{Time-indexed indistinguishability on Phys}}
% ------------------------------------------------------------
Fix a compatible irreversible descent. Thus a
\(G\)-equivariant reversible flow \((\Phi_t)_{t\in\mathbb R}\) on \(X\)
is given together with its descended reversible flow
\((\widehat\Psi_t)_{t\in\mathbb R}\) on \(\Phys\), and an internal
observer reduction
\[
	X\xrightarrow{\pi}\Phys\xrightarrow{\theta}Y,
	\qquad
	\rho=\theta\circ\pi.
\]
\begin{definition}[Kernel filtration on \(\Phys\)]
	For each \(t\ge0\), define an equivalence relation
	\[
		\mathcal F_t:=\ker(\theta\circ\widehat\Psi_t)
		\subseteq \Phys\times \Phys
	\]
	by
	\[
		(p_1,p_2)\in \mathcal F_t
		\iff
		\theta(\widehat\Psi_t(p_1))=\theta(\widehat\Psi_t(p_2)).
	\]
\end{definition}
\begin{lemma}[Equivalence]
	Each \(\mathcal F_t\) is an equivalence relation on \(\Phys\).
\end{lemma}
\begin{proof}
	Fix \(t\ge0\). Since \(\mathcal F_t\) is the kernel relation of the map
	\[
		f_t:=\theta\circ\widehat\Psi_t\colon \Phys\to Y,
	\]
	the relation
	\[
		(p,q)\in\mathcal F_t
		\iff
		f_t(p)=f_t(q)
	\]
	is reflexive, symmetric, and transitive.
\end{proof}
\begin{lemma}[Monotone coarsening]
	\label{lem:flow-Ft-monotone}
	If \(0\le s\le t\), then
	\[
		\mathcal F_s\subseteq\mathcal F_t.
	\]
\end{lemma}
\begin{proof}
	Let \((p_1,p_2)\in\mathcal F_s\). By definition,
	\[
		\theta(\widehat\Psi_s(p_1))
		=
		\theta(\widehat\Psi_s(p_2)).
	\]
	Write
	\[
		t=s+(t-s).
	\]
	Applying \(\Psi_{t-s}\) to both sides gives
	\[
		\Psi_{t-s}\bigl(\theta(\widehat\Psi_s(p_1))\bigr)
		=
		\Psi_{t-s}\bigl(\theta(\widehat\Psi_s(p_2))\bigr).
	\]
	Using \cref{lem:flow-factor-identity},
	\[
		\Psi_{t-s}\circ\theta
		=
		\theta\circ\widehat\Psi_{t-s},
	\]
	and then the flow law on \(\Phys\), we obtain
	\[
		\theta(\widehat\Psi_t(p_1))
		=
		\theta(\widehat\Psi_t(p_2)).
	\]
	Hence \((p_1,p_2)\in\mathcal F_t\), so
	\[
		\mathcal F_s\subseteq\mathcal F_t.
	\]
\end{proof}
\begin{lemma}[Coarsening map between quotient partitions]
	\label{lem:flow-quotient-surjection}
	If \(0\le s\le t\), then the identity map on \(\Phys\) induces a
	canonical surjection
	\[
		q_{s,t}\colon \Phys/\mathcal F_s \twoheadrightarrow \Phys/\mathcal F_t,
		\qquad
		[p]_{\mathcal F_s}\longmapsto [p]_{\mathcal F_t}.
	\]
\end{lemma}
\begin{proof}
	Well-definedness follows from \cref{lem:flow-Ft-monotone}: if
	\[
		[p]_{\mathcal F_s}=[p']_{\mathcal F_s},
	\]
	then
	\[
		(p,p')\in \mathcal F_s\subseteq \mathcal F_t,
	\]
	hence
	\[
		[p]_{\mathcal F_t}=[p']_{\mathcal F_t}.
	\]
	Surjectivity is immediate.
\end{proof}
\begin{remark}[Arrow of time]
	The filtration
	\[
		\mathcal F_0\subseteq \mathcal F_s\subseteq \mathcal F_t
		\qquad
		(0\le s\le t)
	\]
	grows monotonically. Thus indistinguishable states can merge but never
	split. The arrow of time is therefore located entirely in the observer
	reduction \(\theta\), not in the intrinsic quotient \(\pi\).
\end{remark}
% ------------------------------------------------------------
\section{Entropy and extended pseudo-ultrametric}
% ------------------------------------------------------------
\begin{definition}[Observer entropy]
	Assume that
	\[
		|\Phys/\mathcal F_t|<\infty
	\]
	for the times under consideration. Define
	\[
		S_\theta(t):=-\log |\Phys/\mathcal F_t|.
	\]
\end{definition}
\begin{proposition}[Entropy monotonicity]
	Assume that \( |\Phys/\mathcal F_t|<\infty \) for the times under
	consideration. Then \(S_\theta(t)\) is nondecreasing in \(t\).
\end{proposition}
\begin{proof}
	If \(0\le s\le t\), then by \cref{lem:flow-quotient-surjection} there
	exists a surjection
	\[
		q_{s,t}\colon \Phys/\mathcal F_s\twoheadrightarrow \Phys/\mathcal F_t.
	\]
	Hence
	\[
		|\Phys/\mathcal F_t|
		\le
		|\Phys/\mathcal F_s|.
	\]
	Applying the decreasing function \(-\log\) yields
	\[
		S_\theta(s)\le S_\theta(t).
	\]
\end{proof}
\begin{remark}[Observer dependence]
	The entropy functional depends on the reduction \(\theta\). Different
	internal observer families induce different filtrations
	\[
		(\mathcal F_t)_{t\ge0}
	\]
	and therefore different entropy functions. What is intrinsic is not the
	numerical value of \(S_\theta(t)\), but its monotone direction, which
	follows purely from kernel growth.
\end{remark}
\begin{definition}[Observer extended pseudo-ultrametric]
	Define a function
	\[
		d_\theta\colon \Phys\times\Phys\to[0,\infty]
	\]
	by
	\[
		d_\theta(p,q)
		:=
		\inf\{t\ge0:(p,q)\in \mathcal F_t\}.
	\]
	If the displayed set is empty, we use the convention
	\(d_\theta(p,q)=+\infty\).
\end{definition}
\begin{proposition}[Extended pseudo-ultrametric structure]
	The function \(d_\theta\) is an extended pseudo-ultrametric on
	\(\Phys\). That is:
	\begin{enumerate}[label=\textup{(\roman*)}]
		\item
		      \[
			      d_\theta(p,p)=0;
		      \]
		\item
		      \[
			      d_\theta(p,q)=d_\theta(q,p);
		      \]
		\item
		      \[
			      d_\theta(p,r)\le \max\{d_\theta(p,q),d_\theta(q,r)\}.
		      \]
	\end{enumerate}
\end{proposition}
\begin{proof}
	Since
	\[
		(p,p)\in \mathcal F_0,
	\]
	one has
	\[
		d_\theta(p,p)=0.
	\]
	Each relation \(\mathcal F_t\) is symmetric, hence so is \(d_\theta\).
	For the strong triangle inequality, set
	\[
		a:=d_\theta(p,q),
		\qquad
		b:=d_\theta(q,r),
		\qquad
		m:=\max\{a,b\}.
	\]
	If \(m=+\infty\), then
	\[
		d_\theta(p,r)\le \max\{d_\theta(p,q),d_\theta(q,r)\}
	\]
	is immediate. So assume \(m<+\infty\).
	Let \(\varepsilon>0\). By definition of infimum, there exist
	\[
		t_1\le a+\varepsilon,
		\qquad
		t_2\le b+\varepsilon
	\]
	such that
	\[
		(p,q)\in \mathcal F_{t_1},
		\qquad
		(q,r)\in \mathcal F_{t_2}.
	\]
	Set
	\[
		t:=\max\{t_1,t_2\}\le m+\varepsilon.
	\]
	By \cref{lem:flow-Ft-monotone},
	\[
		(p,q)\in \mathcal F_t,
		\qquad
		(q,r)\in \mathcal F_t.
	\]
	Since \(\mathcal F_t\) is an equivalence relation, it is transitive, so
	\[
		(p,r)\in \mathcal F_t.
	\]
	Hence
	\[
		d_\theta(p,r)\le t\le m+\varepsilon.
	\]
	Letting \(\varepsilon\downarrow0\) yields
	\[
		d_\theta(p,r)\le \max\{d_\theta(p,q),d_\theta(q,r)\}.
	\]
\end{proof}
\begin{remark}[Coarsening time as geometry]
	The extended pseudo-ultrametric \(d_\theta\) packages the filtration
	\((\mathcal F_t)_{t\ge0}\) as a geometry on \(\Phys\):
	states are close when they become \(\theta\)-indistinguishable early.
	Pairs that never become \(\theta\)-indistinguishable simply remain at
	infinite distance. This geometry is observer-dependent, but the strong
	triangle inequality is purely filtration-theoretic.
\end{remark}
% ------------------------------------------------------------
\section{The two-locus residue}
% ------------------------------------------------------------
We now return to the enrichment question. The point is that observer
reduction changes distinguishability at the level of reduced dynamics,
but does not change the internal classification of compositional lifts
back to representatives.
\begin{theorem}[Irreversible descent preserves the two-locus residue]
	Under the standing principle \cref{sp:closed-world}, let
	\[
		\rho=\theta\circ\pi\colon X\twoheadrightarrow Y
	\]
	be a compatible irreversible descent.
	Then:
	\begin{enumerate}[label=\textup{(\arabic*)}]
		\item the reduced forward dynamics on \(Y\) is irreversible in the
		      sense that it admits no two-sided extension intertwining \(\rho\);
		\item quotient semantics remains determined by the intrinsic quotient
		      \[
			      \pi\colon X\to\Phys;
		      \]
		\item every compositional lift of quotient-level protocol data to
		      representatives in \(X\) remains exhausted by the same two
		      previously classified loci, possibly both and with no third
		      locus:
		      \begin{enumerate}[label=\textup{(\alph*)}]
			      \item representative sections;
			      \item morphism-level transport/holonomy data.
		      \end{enumerate}
	\end{enumerate}
\end{theorem}
\begin{proof}
	Item \textup{(1)} is exactly \cref{lem:flow-information-loss}:
	information loss precludes any two-sided extension of the reduced
	forward semigroup intertwining \(\rho\).
	For item \textup{(2)}, the quotient semantics of the closed system is
	unchanged by the observer reduction. The intrinsic physical state space
	is \(\Phys=X/G\), obtained from the canonical diagonal quotient. The map
	\[
		\theta\colon \Phys\twoheadrightarrow Y
	\]
	is an additional internal observer reduction imposed \emph{after} that
	quotient. It may erase distinctions between physically distinct states,
	but it does not alter the intrinsic quotient semantics itself.
	For item \textup{(3)}, any compositional enrichment of quotient-level
	protocol data back to representatives in \(X\) is a representative-lift
	problem through the action groupoid. By the lift-classification theorem
	of Chapter~\ref{ch:lift} and the descent-obstruction classification of
	Chapter~\ref{ch:descent_obstruction}, such a lift is exhausted by the
	same two kinds of additional data: object-level representative choice
	and
	morphism-level transport cocycle or holonomy data. Since the observer
	reduction \(\theta\) acts only after passage to \(\Phys\), it creates no
	new action-groupoid lift data above quotient semantics. Hence no third
	enrichment locus appears.
\end{proof}
\begin{remark}[Structural summary]
	The logical architecture is therefore rigid:
	\[
		X
		\xrightarrow{\ \pi\ }
		\Phys
		\xrightarrow{\ \theta\ }
		Y.
	\]
	The first map is intrinsic and reversible; it defines the closed-system
	quotient semantics. The second map is internal but observer-dependent;
	it may be irreversible and induces the monotone indistinguishability
	filtration on \(\Phys\). Yet even in the irreversible regime,
	representative choice and morphism-level transport or holonomy data
	remain the only enrichment loci above quotient semantics, possibly both,
	with no third locus, by
	\cref{thm:lift-classification-main,cor:two-loci-lift}.
\end{remark}
% ------------------------------------------------------------
\section{Interpretation}
% ------------------------------------------------------------
The chapter establishes four points which should be kept sharply
separate.
First, the canonical quotient
\[
	\pi\colon X\to\Phys
\]
is not an irreversible coarse-graining. It removes only the intrinsic
diagonal redundancy determined by closed comparison semantics. Because
\(G\)-equivariant reversible flows descend through \(\pi\) as
reversible flows, no arrow of time originates there.
Second, observer reduction is internal to the closed system. It is not
an externally imposed surjection, but the quotient induced by a proper
restricted family of admissible quotient-level observations. Thus the
second descent stage
\[
	\theta\colon \Phys\twoheadrightarrow Y
\]
arises from observer limitation already encoded inside the comparison
stack.
Third, once a reversible flow on \(X\) and its descended reversible
flow on \(\Phys\) are fixed, irreversibility appears only when such an
internal observer reduction is forward-compatible with that descended
flow but not backward-stable. The exact mechanism is backward fiber
separation. This induces the monotone kernel
filtration
\[
	(\mathcal F_t)_{t\ge0}
\]
on \(\Phys\), from which observer entropy and the extended
pseudo-ultrametric geometry follow canonically.
Fourth, the enrichment residue is unaffected by this observer-level
irreversibility. Even when reduced dynamics becomes irreversible on
\(Y\), the only compositional ways to lift quotient-level data back to
\(X\) remain the two loci of
\cref{thm:lift-classification-main,cor:two-loci-lift}: representative
choice and morphism-level transport or holonomy data, possibly both,
with no third enrichment locus.
In particular, the chapter introduces no new primitive for irreversible
behavior. Irreversibility is carried by internal observer-induced
failure of descent, not by an enlargement of the intrinsic comparison
stack.
\begin{remark}[What is not yet used]
	Nothing in the present chapter uses smooth realization, connection
	theory, curvature, graded commutator structure, or quadratic carriers.
	The role of the chapter is more basic: for fixed reversible flow data,
	it isolates the exact observer-level locus at which irreversibility can
	appear once a proper internal observer family is forward-compatible, and
	it proves that representative choice and morphism-level transport or
	holonomy data remain the only enrichment loci above quotient semantics,
	possibly both, with no third locus. \Cref{ch:stitching} is the
	immediate next continuation, and \cref{ch:causality,ch:connection} use
	the stabilized arena later, but none of that is used here.
\end{remark}

\section{Conclusion}
\label{sec:flow-conclusion}
The chapter isolates irreversibility at its exact structural location:
not in the intrinsic quotient map \(X\to\Phys\), but in the subsequent
observer-induced descent \(\Phys\twoheadrightarrow Y\) once fixed
reversible flow data are in place and a proper internal observer family
is forward-compatible with the descended flow but fails to be
backward-stable. The arrow of time is therefore derived internally
rather than postulated externally.

Accordingly, \cref{ch:stitching} stabilizes this dynamical picture by
assembling refinement and observer-time quotients into a canonical
two-parameter inverse-limit arena later used in
\cref{ch:causality,ch:connection}.